This section reads the phenomenon as a security risk under the Artificial Intelligence Risk Management Framework (AI RMF) of the National Institute of Standards and Technology (NIST).

The forecaster sits in a control loop in which its output drives battery dispatch and curtailment on a photovoltaic micro-grid, so a forecasting failure can become a control failure, and the property tested here is a security property as well as an accuracy one.

Phase-locking is deterministic, attacker-independent and present at deployment time, and the candidate set follows from $(f_s,P,S)$ alone, marking where a blind spot is most expected without excluding others. Where an exposure would sit, and how it moves with the configuration, is known before anything is measured: enough to place a monitor and to bound what a mitigation could reach, not enough to claim that a loss occurs.

\subsection{Risk Management Framework}
The AI RMF\autocite{tabassiArtificialIntelligenceRisk2023} is the reference for that reading. It separates the trustworthiness characteristics at stake, validity and reliability, security and resilience, from the governance functions that turn evidence into a control, and asks that known limitations be mapped, measured and managed rather than merely disclosed.

Its four functions assign responsibility and oversight (GOVERN), fix the deployment context and the pathways to harm (MAP), track the likelihood, severity and uncertainty of each risk (MEASURE), and select and monitor responses (MANAGE). They inform one another continuously, so \autoref{tab:riskRegister} records the functions each threat engages rather than a sequence.

None of the evidence each pathway needs establishes a loss (\autoref{tab:bayesEstimates}): before the gates the forecast recovers a tone better at a candidate than beside it, the representational cost lies within the equivalence band, and the stride sites sit where $f_s/S$ predicts. Every reading fails the gates, so the register records them as the direction of the evidence, lowers no residual on their strength, and keeps open every pathway that depends on a loss.

Each row of \autoref{tab:riskRegister} names a pathway, the functions it engages, the estimand that would establish it, the available control and the residual exposure, which stays conditional until that estimand passes the gates.

\begin{table*}[p]
    \centering
\small
    \setlength{\tabcolsep}{4.5pt}
    \begin{tabular*}{\textwidth}{@{\extracolsep{\fill}}>{\raggedright\arraybackslash}m{0.45cm}>{\raggedright\arraybackslash}m{3.35cm}>{\raggedright\arraybackslash}m{1.70cm}>{\raggedright\arraybackslash}m{3.85cm}>{\raggedright\arraybackslash}m{3.25cm}>{\centering\arraybackslash}m{1.45cm}@{}}
        \toprule
        \textbf{ID} & \textbf{Risk} & \textbf{AI RMF} & \textbf{Evidence in this report} & \textbf{Control available} & \textbf{Residual} \\
        \midrule
        $T1$ & Adversarial masking:\newline a periodic component injected on a candidate frequency may be under-represented and escape forecast-residual detection & \shortstack[l]{MAP\\MEASURE} & Precondition is the support rule $e^{\theta_{\mathrm{lock}}}>1.2$ (Model~B) with the site locations of $\kappa_S$ (Model~D2); the pathway itself needs the paired recovery ratio of Model~A. Pre-gate: $e^{\theta_{\mathrm{lock}}}=1.06$, $e^{\bar\beta}=1.35$, a gain; not reportable & Spectral monitor on the raw stream, upstream of tokenisation, at the enumerated candidate set & high \\
        \cmidrule(lr){1-6}
        $T2$ & Systematic blind spot with no adversary:\newline a legitimate plant harmonic sits on a candidate frequency and may be persistently under-represented & \shortstack[l]{MAP\\MEASURE\\MANAGE} & Candidate set is closed-form in $(f_s,P,S)$ and needs no fit; the cost at those sites is decided by $e^{\theta_{\mathrm{lock}}}$. Pre-gate: $1.06$, inside the equivalence band; not reportable & Choose $(P,S)$ against the known spectral content of the plant; \valueof{PartiallyObservable} degradation & medium \\
        \cmidrule(lr){1-6}
        $T3$ & Configuration drift:\newline a change of sampling interval, patch size or checkpoint geometry silently relocates the candidate set & \shortstack[l]{GOVERN\\MAP} & $H3a$: relocation is predicted in closed form by $f_s/S$, and whether the sites follow it is decided by $\kappa_S$. Pre-gate: $\kappa_S=1.000$ $[0.990,1.011]$; not reportable. The residual rests on the closed-form set, not on this reading & Recompute the candidate set from the knowledge base at every configuration change & low \\
        \cmidrule(lr){1-6}
        $T4$ & False assurance: overlap is increased and treated as a fix & \shortstack[l]{GOVERN\\MANAGE} & $M1$ is decided by $\delta_O$ under a two-leg rule; pre-gate both legs fail, $\delta_O=-0.13$, inconclusive and not reportable. Independently of any fit, the patch branch is invariant to $S$ by construction & Do not credit overlap as a control; record the unestablished status in the model card & high \\
        \cmidrule(lr){1-6}
        $T5$ & Defence placed inside the blind spot: phase randomisation, or anomaly detection on forecast residuals & MANAGE & $H2$ is decided by $\sigma_\phi$ (Model~C): a deficit independent of phase makes randomising it useless. Pre-gate: $\sigma_\phi=0.006$; not reportable & Move detection upstream of tokenisation; treat residual-based detection as insufficient alone & medium \\
        \bottomrule
    \end{tabular*}
\caption{Risk register. AI RMF names the function of the framework\autocite{tabassiArtificialIntelligenceRisk2023} engaged by each item. Residual is the exposure left by the control, judged qualitatively: high where no operator control closes it, medium where the control is partial or needs design-time choices, low where a procedure that follows from the arithmetic closes it. No residual is lowered on estimates that fail the gates.}
    \label{tab:riskRegister}
\end{table*}


\subsection{Adversarial-masking threat model}
A weakness becomes a vulnerability when someone can reach it. Because phase-locking is deterministic, an adversary needs no model access, query budget or knowledge of the training corpus, only the configuration $(f_s,P,S)$ of the deployed forecaster, from which $\mathcal F_{\mathrm{lock}}$ follows by arithmetic. In the taxonomy of adversarial machine learning\autocite{vassilevAdversarialMachineLearning2024} the weakness is an availability violation, since what the model would lose at a candidate frequency would be lost for every input and user. It becomes an integrity violation of the evasion type once an adversary uses that loss to hide a manipulation.\footnote{Deliverable~2 classified the weakness as an availability violation only. The classification is extended because the masking threat model developed here seeks to hide a manipulation rather than to degrade the service, which the taxonomy treats as evasion.}

\autoref{tab:threatModel} sets out the adversarial-masking threat model in five elements, from the actor's capability and access path to the manipulation, the consequence and the exposure left after the mitigation. Every row is conditional: it describes the exposure that would follow if a loss at the candidate sites were confirmed, not one that has been observed.

\begin{table}[!t]
    \centering
\small
    \setlength{\tabcolsep}{5pt}
    \begin{tabular}{@{}M{0.215\linewidth} M{0.735\linewidth}@{}}
        \toprule
        \textbf{Element} & \textbf{Reading for this system} \\
        \midrule
        Adversary capabilities & No model access of any kind and no knowledge of the training corpus, only knowledge of the configuration $(f_s,P,S)$ and the ability to drive a periodic load on the photovoltaic plant, which an inverter switching schedule or a switchable load already provides. \\
        \addlinespace[3pt]
        Attack surface & The raw telemetry stream before tokenisation, reached by acting on the physical process or on the data path. \\
        \addlinespace[3pt]
        Manipulation of telemetry & A small periodic component at a candidate frequency, which is legitimate plant behaviour and survives any plausibility check on the raw signal. \\
        \addlinespace[3pt]
        Operational consequence & The component is present in the plant and, if structural aliasing is confirmed, under-represented in the model, so dispatch and curtailment are decided on an incomplete picture, and a detector reading the forecast residual is looking through the same tokeniser that lost it. \\
        \addlinespace[3pt]
        Residual risk & The pre-registered mitigation raises overlap by shortening the stride at fixed patch size, so it thins the stride branch alone and $\mathcal{F}_{P}$ survives at every deployed configuration. Shrinking the patch instead thins $\mathcal{F}_{P}$, and with it $\mathcal{F}_{S}$ once $S \le P$ forces the stride down as well; but that raises both fundamentals rather than closing either comb, so the candidate set is relocated and thinned, not removed. \\
        \bottomrule
    \end{tabular}
\caption{The adversarial-masking threat model in five elements, from the capability the actor needs to the exposure left after the pre-registered mitigation. The first row sets it apart: what is required is not access to the model but the ability to make the plant do something it is entitled to do, so nothing in the manipulation is anomalous where a plausibility check would sit.}
    \label{tab:threatModel}
\end{table}


This is the adversarial-masking row of \autoref{tab:riskRegister}, carried at high residual: the injected component is legitimate plant behaviour, which a monitor on the raw stream can locate at a candidate frequency but cannot tell from ordinary operation.

\subsection{Threat models beyond the adversarial one}

The framework asks for systematic vulnerability, not only for attacks, and the more likely exposure has no adversary. On one-minute telemetry, \autoref{eq:branches} puts the patch branch of a $P{=}16$ tokeniser at periods of $16/k$ minutes: $16$, $8$, $5.3$, $4$, $3.2$, $2.7$ and $2.3$ minutes. These are the timescales of inverter duty cycling and cloud-shadow transients, ordinary plant behaviour sitting on the frequencies at which the model may be least able to represent it. If the loss were confirmed, a monitor built on such a forecaster would be systematically less sensitive there, for reasons no operational log would reveal. This is the systematic blind spot of \autoref{tab:riskRegister}, and the reason the candidate set belongs in the deployment documentation.

The third exposure, configuration drift, is a matter of governance. Because the candidate set is a function of configuration, changes that no one would route through a security review relocate it: resampling telemetry from one minute to thirty seconds, swapping to a checkpoint of another geometry, or choosing a patch size for latency. $H3a$ would make this tractable: if the sites follow $f_s/S$, recomputing the candidate set after a change is arithmetic rather than re-measurement.

\subsection{Controls and residual exposure}

Two controls in this design are real. The first answers the defence placed inside the blind spot ($T5$ in \autoref{tab:riskRegister}): detection has to sit before tokenisation, on the raw telemetry, since anything placed after it would inherit any blind spot the tokeniser induces. A spectral monitor on the raw stream is inexpensive because the candidate set is closed-form. The second is the semantic layer of \autoref{sec:appOntology}, which degrades actuation rather than failing silently: a \valueof{PartiallyObservable} assessment reduces the battery cap and lowers decision confidence, so uncertainty about the representation reaches the actuator instead of being resolved by assumption. This is the resilience the framework asks for, and it does not depend on how the behavioural claim resolves.

One control cannot be credited. The pre-registered mitigation, increased patch overlap, is not established: both legs of its rule fail and the reading is not reportable. Even if established it would be partial, since overlap moves the stride branch and leaves the patch branch untouched. An operator who raises the overlap and records the risk as mitigated has bought the false assurance of \autoref{tab:riskRegister}, which is why it is carried at high residual.
